\documentclass{article}

\usepackage{sectsty}
\usepackage{graphicx}
\usepackage{csquotes}
\usepackage[
backend=biber,
style=alphabetic,
sorting=ynt
]{biblatex}
\addbibresource{biblio.bib} %Import the bibliography file

\usepackage{url}
\usepackage{hyperref}

% Margins
\topmargin=-0.45in
\evensidemargin=0in
\oddsidemargin=0in
\textwidth=6.5in
\textheight=9.0in
\headsep=0.25in

\title{Comparative review}

\begin{document}
\maketitle	

\paragraph{Allgor2023 \cite{Allgor2023}}
L'articolo si concentra sulla determinazione di una picking policy, senza affrontare i problemi di stowing o storage, introducendo vincoli sul soddisfacimento dei target rate per ciascun process path, un upper bound sul numero di ordini parzialmente prelevati e una allocazione dei pick coerente con la capacità produttiva delle diverse zone per evitare innecessario idle time. 

Il modello matematico si divide in un \emph{picking module} e un \emph{assignment module}. Il primo modulo considera gli ordini da prelevare e le loro due-date, l'inventorio, la capacità lavorativa, il throughput, e i target dei cycle-time e determina la pick-window a l'assegnazione oridini-pod con l'obiettivo di massimizzare il pile-on. Il secondo prende come input la soluzione del primo modulo e assegna gli ordini alle working station minimizzando la distanza percorsa dai robot e l'idle time. 

Il primo modulo è di fatto intrattabile computazionalmente e viene quindi risolto mediante euristica che suddivide in 4 sottoproblemi la decisione. 

L'articolo fa un'analisi di come la size della pick-window impatti sul cycle-time di ordini con più SKU e sul pile-on.




\paragraph{Barnhart2024 \cite{Barnhart2024}}
L'articolo formula un problema denotato come TDS-CW, Task Design and Scheduling problem with Congestion and restricted Workload; il modello massimizza il throughput (ossia, massimizza il numero di oggetti presi nel planning horizon) introducendo vincoli sul numero di ordini massimo non completati e sulla congestione (questi sono formulati con una time-space network). 

Dal momento che il problema così definito è intrattabile per istanze medio-grandi, si risolve il problema con un algoritmo di ricerca locale in cui i vicinati vengono generati tramite un algoritmo di machine learning. Nello specifico, in una fase preliminare di training si determina una funzione capace di predire il valore della funzione obiettivo e poi grazie a questa si genera un sottoproblema in cui si masssimizza il miglioramento predetto da questa.



\paragraph{DeKoster2023 \cite{DeKoster2023}}
L'articolo si pone il problema della schedulazione di un batch fissato di ordini senza considerare vincoli di congestione e e propone due approcci differenti. 

Il primo si basa su una decomposizione in tre sottoproblemi successivi (clustering degli ordini e assegnazione cluster-pod, assegnamento alle workstation, assegnamento robot-pod) con obiettivi di minimizzazione delle distanze percorse. Il secondo è una formulazione integrata che decide congiuntamente assegnamenti e percorsi, con l'obiettivo di minimizzare il makespan del batch.
Per risolvere la seconda fase viene introdota una metaeuristica, \enquote{Two-Layer Revolving Algorithm}, in cui viene iterativamente risolto il modello (o una sua versione rilassata) mantenendo alcune variabili decisionali fissate in modo da accellerare la ricerca dell'ottimo.

La trattazione include inoltre analisi di sensitività dell'impatto che la composizione degli ordini, il layout del FC e la posizione delle SKU nei pod hanno sul makespan.
Infine, si cerca di superare la limitazione introdotta dal \enquote{batch-mode} ad un caso in cui al tempo 0 ci siano già ordini in lavorazione.



\paragraph{Merschformann2019 \cite{Merschformann2019}}
L'articolo si concentra sul definire e testare delle policy per i seguenti problemi decisionali posti a livello operazionale: \emph{Pick Order Assignment} (selezione di un pick order dal backlog e assegnameno ad una workstation), \emph{Replenishment Order Assignment} (assegnazione del task di replenishment, che vengono serviti con politica FCFS, ad una replenishment workstation), \emph{Pick Pod Selection} (selezione di un pod per un task di tipo pick), \emph{Replenishment Pod Selection} (selezione di un pod per un task di tipo replenishment), \emph{Pod Storage Assignment} (scelta di dove far ritornare il pod).

Per valutare le regole si è usato un software di simulazione agent-based e event-driven e si è proceduto in due fasi.
Inizialmente, si è fissato il layout del FC e si sono valutate singolarmente le regole per ogni problema in base ad alcune misure di prestazione (che si sono tutte rivelate essere fortemente dipendenti dal throughtput). In seguito, si è valutato come le policy che si erano rivelate promettenti si sarebbero comportate in differenti layout e scenari di domanda.


\paragraph{Boysen2023 \cite{Boysen2023}}
Lo scenario in cui si pone l'articolo è quello in cui le SKU arrivano sequenzialmente e il compito dei robot è di prendere ciascuna SKU e portarla al suo collection point. In questo senso il problema si divide in 3 decisioni: Piece-to-Order assignment (P2O), Order-to-Collection Point assignment (O2CP) and Piece-to-Robot assignment and robot scheduling (P2R).

L'articolo prima considera il caso deterministico di informazione perfetta e risolve il modello, che vorrebbe minimizzare il makespan,  decomponendolo nelle sue tre decisioni: P2O (minimizzazione del makespan), O2CP (minimizzazione della distanza percorsa) e P2R. Questi vengono risolti tramite euristiche basate su regole di priorità.

Si estende la trattazione al caso di informazione parziale (\emph{limited lookahed}): ora, per pendere una decisione, vengono prima generati degli scenari, si risolvono i corrispondenti problemi deterministici e, infine, la decisione è presa in base ad una funzione di consenso.

L'articolo si chiude con un'analisi di sensitività sull'impatto che avrebbe l'aumentare il numero di robot o la capacità di lookahed.


\paragraph{Xu2024 \cite{Wu2024}}
L'articolo cerca di risolvere in maniera integrata sia la gestione degli ordini ma anche il task di replenishment. La performance del sistema è misurata in base al makespan.

Per fare ciò si suddivide il problema in 5 fasi: (1) si suddividono gli ordini in cluster e si assegnano questi alle picking station; (2) si assegnano i pod agli ordini cercando di minimizzare i replenishment necessari; (3) si determinano i percorsi dei pod minimizzando la massima distanza percorsa; (4) si assegnano i robot ai pod minimizzando la distanza del robot dal pod; (5) schedulazione delle attività dei pod minimizzando il makespan. Ogni fase è modellata tramite MILP.

I modelli vengono risoltri tramite Variable Neighborhood Search ponendosi il problema di come generare una buona soluzione iniziale (risolto tramite euristica) e come definire un buon vicinato (vengono definiti 4 operatori).

L'articolo si chiude con un'analisi di sensitività sulla composizione degli ordini e sul numero di picking station.



\paragraph{Subramanian2024 \cite{Subramanian2024}}
L'articolo presenta l'algoritmo SASORL (Simultaneous Allocation and Sequencing of Orders Reinforcement Learning) che simultaneamente alloca e sequenza gli ordini ai robot cercando di minimizzare la distanza percorsa. In questo articolo, con il termone \enquote{ordine} non ci si riferisce alla classica spedizione, ma ad uno scaffale da spostare.

Il framework proposto è un algoritmo di reinforcement learning in cui lo stato è rappresentato dagli ordini soddisfatti, mentre lo spazio delle azioni sono gli ordini da soddisfarre. Per costruire la Q-table, si considera un termine di penalizzazione basato sulla distanza che bisogna percorrere per soddisfarre l'azione scelta. 
Ad ogni iterazione l'algoritmo impara come minimizzare il termine di penalità.



\newpage
\paragraph{Shi2024 \cite{shi2024}}
L'articolo definisce una strategia congiunta di ottimizzazione dell'assegnamento degli ordini e quello dei pod senza porsi il problema dedlla possibile congestione dei robot o di un eventuale replenishment.
Il modello matematico formulato è un MIP, che vuole minimizzare il numero totale di visite dei pod alle workstation, di fatto intrattabile per istanze medio-grandi. 

Per risolvere il modello viene proposto un'euristica, denotata come FFSA (Fulfillment-Focused Simultaneous Assignment), che si sviluppa in due fasi.
Nella prima fase, viene usato un hybrid ALNS per determinare un sottoinsieme di pod per soddisfare la domanda in cui gli operatori di distruzione sono definiti classicamente, mentre l'operatore di costruzione si ottiene risolvendo un problema di ottimizzazione. Nella seconda fase, si definisce l'euristica MRACS che si basa su una prima strategia iterativa per determinare un iniziale assegnamento pod-(alcuni) ordini e poi di uno schema MRB per ottenere l'intero assegnamento (strategia construttiva basata sul massimo marginal return ottenibile).

Nell'articolo c'è una descrizione su come generare istanze sintetiche.









\newpage
\printbibliography
\end{document}